% !TEX root = ./main.tex

In a quest for a notion of adaptation, the authors of~\cite{bcglv12}
advocate a conceptual framework grounded on a neat separation of
concerns.
%
Starting from the observation that the behaviour of a component is
determined by a program defining \emph{control} and manipulating
\emph{data}, adaptation requires to consider, and possibly, modify
specific computational or physical resources.
%
The conceptual framework in~\cite{bcglv12} envisages \emph{adaptation}
as the possibility of a component to manipulate and take decisions
according to a well-identified set of \emph{control data}.
%
For instance, an adaptive behaviour may need to account for memory
usage, network traffic, or environmental physical conditions like
e.g., temperature.

We render the conceptual framework in our platform as detailed in
\cref{sec:arch,sec:control-loop}.

\subsection{An architecture for runtime adaptation}\label{sec:arch}
%
Embracing the conceptual framework introduced in~\cite{bcglv12} allows
one to define adaptation as the property of a component of controlling
and possibly modifying at run-time its behaviour based on a specific
set of control data.

This acceptation implicitly requires a \emph{monitor} that can collect
and make available the control data necessary for enforcing adaptive
behaviour.
%
Noteworthy, the nature of control data impacts on the complexity of
the monitoring activities.
%
When control data are \emph{exogenous}, namely they originate outside
the application (e.g., from sensors, operating system, etc.), their
usage in the adaption process can be simply attained by requesting the
information to the monitor (e.g., with a request-response
interaction).
%
Otherwise, for \emph{endogenous} control data, namely data produced by
some components of the application, monitor and application are more
coupled and their interactions can require the definition of specific
application-level protocols.

Leveraging the API of \SEArch, we propose a programming pattern to
support the conceptual framework proposed
in~\cite{bcglv12}.
%
We focus on exogenous control data, leaving adaptation endogenous
control data for future work.
%

\begin{figure}[t!]\centering
  \includegraphics[width = \textwidth]{adaptable-service}
  \caption{Architectural view of our programming pattern for adaptation}
  \label{fig:adaptable_service}
\end{figure}
%
\cref{fig:adaptable_service} provides a high-level architectural view
of an adaptable agent implemented over our platform; the two main
elements of the architecture are the \emph{adaptable service} itself
and the local \emph{middleware} (dark blue box) that sets
up the communication infrastructure necessary for the dynamic composition
of services.
%
As shown, the adaptable service runs behind the middleware,
which mediates the interactions of the services/applications in
execution on a node with every other component in the network,
including other middlewares and the service broker.
%
The mediation basically connects (via the channels
$\textsf{CC}_\textsf{1}, \cdots, \textsf{CC}_\textsf{n}$) the required
ports of the adaptable service with the provided ports of brokered
external components.
%


The structure of an adaptable service is rather complex and consists
of:
\begin{itemize}
\item a collection of contracts (specified as extended CFSMs)
  corresponding to each of the execution scenarios
  ($\{\mathsf{Communication\ contract\ i}\}_\mathsf{1 \leq i \leq n}$
  in the topmost part of the orange box, so that each
  $\mathsf{Communication\ contract\ i}$ corresponds to the
  middleware's channel $\mathsf{CC}_\textsf{i}$; we also emphasise
  this correspondence with colours),
\item a collection of implementations corresponding to each of the
  execution scenarios
  $\{\mathsf{Execution\ scenario\ i}\}_\mathsf{1 \leq i \leq n}$,
  each equipped with an \emph{adaptation contract}
  $\mathcal{E}_\mathsf{i}$ formalising the
  environmental conditions under which the execution scenario should
  be run (bottom-most part of the inner orange box), and
\item a decision procedure that runs a control loop, which uses the control data received from the monitor to determine an appropriate execution scenario to run (leftmost part of the orange box).
\end{itemize}

We remark that each channel $\textsf{CC}_\textsf{i}$ is used by the
middleware to trigger the adaptation scenarios according to the
$\mathsf{Communication\ contract\ i}$.
%
The details of the control loop and of the adaptation contracts of the
execution scenarios are given in \cref{sec:control-loop}.


\subsection{An implementation of the control loop}\label{sec:control-loop}
%
As discussed above, taking inspiration from~\cite{bcglv12}, we
understand adaptability as the capacity of an agent to respond to
changes in the environment.
%
Thus, a significant aspect of our approach relies on the exogenous
control data received from the monitor, and how the decision making
process uses it to determine an appropriate execution scenario at
runtime.

Our programming pattern relies on a fixed set $\mathcal{A}$ of
(\emph{quantitative}) \emph{attributes} (ranged over by $a$) which
store the exogenous control data.

\begin{example}\label{exe:attributes}
  For the sake of our running example, we make the following assumptions:
  \begin{itemize}
  \item \srv{Local DB} and \srv{WebUI} run on the same node while the image
	 provider \srv{Image provider} runs on a different node (cf.
	 \cref{fig:external-service-arch} and \cref{fig:fallback-contract-arch} in \cref{sec:teapat});
  \item adaptations are triggered depending on some constraints
	 involving the attributes $L$ (network latency), $R$
	 (responsiveness of \srv{Image provider}), and $C$ (CPU load at the node of \srv{WebUI}).
  \end{itemize}
  The adaptation contracts of our running examples will be constraints
  on these attributes.
  \finex
\end{example}

The programming pattern proposed here essentially hinges on \SEArch
features to instantiate communication channels used in the control
loop to manage adaptation.
%
More precisely,
%
\begin{itemize}
\item the invocation of a service triggers the
  registration of the $\mathsf{n}$ alternative communication channels in the 
 middleware of the  local node of the service,
\item the control loop managing the runtime adaptations consists of
  iterating the following steps:
  \begin{enumerate}[a.]
  \item acquire from the most recent control data from the monitor (cf. \cite{bcglv12})
	 as an assignment of values to the attributes in $\mathcal A$;
  \item\label{it:select} select a suitable execution scenario
	 according to the control data provided by the monitor;
  \item trigger the execution of the code implementing the selected scenario;
  \item stop the running execution scenario if the updated control data 
          determine that the current execution scenario is no longer admissible
	 (i.e., the contract of the running scenario is invalidated by
	 newly acquired control data) and repeat from step~\ref{it:select}.
  \end{enumerate}
\end{itemize}

The control loop sketched above is designed to select one of the
possible execution scenarios according to the control data provided by
the monitor.
%
More precisely, the identification of a suitable scenario depends on the
control data returned by the monitor, which we render as
mappings $m : \mathcal{A} \to \Re$ that assign a real value to each
attribute in $\mathcal A$.
%
%\begin{example}\label{exe:cd}
%\end{example}
%
This design decision is supported by the framework introduced
in~\cite{lopezpombo:tcs-1043,lopezpombo:ictac23} which guaranties the
existence of a decision procedure for determining satisfiability of a
class of constraints expressed in a first-order logic with
equality whose formulae predicate on real values.\footnote{%
  The framework is based on \emph{real closed fields}~\cite{tarski51},
  a decidable first-order theory of the real
  numbers~\cite[Theorem~37]{tarski51}) to specify contracts on
  quantitative aspects of software artefacts.
  %
  For the sake of this paper, it is enough to assume that contracts
  consist of decidable constraints; the interested reader is referred
  to~\cite{lopezpombo:tcs-1043,lopezpombo:ictac23} and references
  therein for details.
  %
}
%
Hence, this framework enables the use of \emph{Satisfiability Modulo
  Theories solvers}~\cite{biere21} (SMT solvers for short) such as
Z3~\cite{demoura:tacas08}, Yices~\cite{dutertre:cav14}, or
CVC5~\cite{barbosa:tacas22}.
% 
% Hereafter, we let $m : a \mapsto v_a$ with $a \in \mathcal{A}$ range
% over control data mappings and
In fact, we formalise adaptation contracts as constraints on real
values.
%
\begin{example}\label{exe:contract}
We set the following conditions\footnote{%
  The adaptation contracts in our running example are pretty basic,
  for simplicity.
  %
  We remark that our platform can handle non-trivial constraints
  expressed as formulae of the decidable first-order logic over real
  closed fields.
} on $L$, $R$, and $C$:
\begin{align*}
  \mathcal{E}_1 & = L < 5ms \quad \land \quad R < 1s
  \\
  \mathcal{E}_2 & = \neg \mathcal{E}_1 \quad \land \quad 0 \leq C \leq .9
  \\
  \mathcal{E}_3 & = \neg \mathcal{E}_1 \quad \land \quad C > .9
\end{align*}
%
The intuition behind these conditions is that the interaction between
\srv W and \srv P is satisfactory when $\mathcal{E}_1$ holds otherwise
an adaptation should take place.
%
We conceive two adaptations when the remote interaction degrades; \srv
W retrieves the images through \srv S when the local CPU is not
overloaded ($\mathcal{E}_2$), otherwise a default image is displayed
($\mathcal{E}_3$).
\finex
\end{example}


% Besides expressing adaptation contracts (attained using the framework in~\cite{lopezpombo:tcs-1043,lopezpombo:ictac23}),
The dynamic selection of execution scenarios appeals to the use of SMT
solvers to decide the satisfiability of the adaptation contracts
associated to the execution scenarios.
%
To tackle this issue, we equip our platform
with the procedure expressed in Pythonic pseudo-code in
\cref{lst:control-loop} which relies on threads and on the following
assumptions on the execution scenarios (cf.
lines~\ref{lst:threads}-\ref{lst:scenario-def-end} of \cref{lst:control-loop}):
\begin{description}
\item[scenarios] the functions implementing each execution scenario
  are collected in a module called \lstinline[style=pyinl]{ExcScenarios} which also
  exports the list \lstinline[style=pyinl]{ES} of such functions; for simplicity,
  we assume that these functions do not take parameters;
\item[coordination] the module \lstinline[style=pyinl]{ExcScenarios} also
  exports a boolean variable \lstinline[style=pyinl]{done}, shared with the control loop, 
  and used to determine when the scenario has completed its execution;
\item[termination] each execution scenario
  is guaranteed to
   atomically set\footnote{We use the syntax \lstinline[style=pyinl]{<<...>>}
	 to mark critical sections.} \lstinline[style=pyinl]{done} to \lstinline[style=pyinl]{True} before
  returning and to terminate.
\end{description}

\begin{figure}[t]
  \centering
\begin{lstlisting}[style=PythonStyle,caption={Pythonic pseudo-code implementing the control loop.},label={lst:control-loop}]
from threading import Thread $\label{lst:threads}$
from ExcScenarios import ES, done $\label{lst:scenario-def-end}$

done = False

while << not done >>:            	        # atomic access $\label{lst:scenario-selection-start}$
  stop, scenario = False, None 
  i = 0
$\label{scenario-selection-start}$  while scenario == None:
    # get monitored data
$\label{monitor-update-selection}$    m = monitor_update() 
$\label{adaptation-condition-check-selection}$    if not SMT(not $\mathcal{E}_i[a \mapsto \mathsf{m}(a)]_{a \in \mathcal{A}}$) == $\ $ SAT:
      # spawn the chosen scenario ES[i]
      scenario = Thread(target = ES[i]) $\label{lst:thread-creation}$
$\label{selection-execution}$      scenario.start() $\label{lsg:thread-start}$
    else:
$\label{scenario-selection-end}$      i = (i $+$ 1) mod len(ES)

  while not (done or stop): $\label{lst:running}$
    old_m = m
    # get monitored data
$\label{monitor-update}$    m = monitor_update()                
    if not old_m == m:
      stop = SMT(not $\mathcal{E}_i[a \mapsto \mathsf{m}(a)]_{a \in \mathcal{A}}$) == $\  $ SAT
    if stop:
      scenario.terminate() $\label{lst:finishing}$
\end{lstlisting}
\end{figure}

The control loop starts by setting\footnote{%
  Notice that there is no need to make an atomic assignment since the
  control loop has not spawned any other thread yet.
%
} the \lstinline[style=pyinl]{done}
flag to \lstinline[style=pyinl]{False}
and entering the loop on line~\ref{lst:scenario-selection-start} which,
until the \lstinline[style=pyinl]{done} flag is not set:
\begin{enumerate}[a.]
\item chooses a scenario to run if none is running
  (lines~\ref{lst:scenario-selection-start}-\ref{scenario-selection-end})
  and
\item keeps checking if a running scenario should be stopped or it
  gracefully finished (lines~\ref{lst:running}-\ref{lst:finishing}).
\end{enumerate}
% 
To select the execution scenario, the control loop iterates over
the list of function \lstinline[style=pyinl]{ES} and stops
as soon as it finds a function whose adaptation contract is true.
%
This is checked by querying an SMT solver whether the negation of
$\mathcal{E}_i[a \mapsto \mathsf{m}(a)]_{a \in \mathcal{A}}$ is
satisfiable or not;\footnote{%
  As usual, $\mathcal{E}[a \mapsto v_a]_{a \in \mathcal{A}}$ denotes the
  application to the formula $\mathcal{E}$ of the substitution
  $[a \mapsto v_a]_{a \in \mathcal{A}}$ which maps each
  attribute $a \in \mathcal{A}$ to a value $v_a$.
  %
}
\begin{itemize}
\item gets the control data from the monitor (line~\ref{monitor-update-selection};
  assume that \lstinline[style=pyinl]{monitor_update} takes
  care of communicating with the monitor and returning the mapping
  to the control loop),
\item if the adaptation contract $\mathcal{E}_i$
  is violated (namely its negation is satisfiable)
  under the returned control data, then the scenario cannot be chosen,
\item otherwise, the scenario is chosen and a thread running the associated implementation
  is started (lines~\ref{lst:thread-creation}-\ref{lsg:thread-start})
  and the execution continues on line~\ref{lst:running}.
\item the loop on lines~\ref{lst:running}-\ref{lst:finishing} keeps
  checking that the adaptation contract of the running scenario is
  still valid; if the contract gets violated, the scenario is
  terminated (line~\ref{lst:finishing}).
\end{itemize}

For simplicity, the thread running the execution scenario is
terminated; we indeed opted for a lighter presentation instead of
devising a more general mechanism to coordinate the thread and the
control loop for a graceful termination (where e.g., the thread
restores some application-dependent invariants before terminating) at
the cost of complicating the presentation.


%%% Local Variables:
%%% mode: LaTeX
%%% TeX-master: "main"
%%% End:
